% !TEX root = main.tex
% !TEX program = xelatex
%%% =========================================================================
%%% Starter Template for ECE Ph.D. Proposal
%%% =========================================================================

\documentclass{eceproposal}

% --- Document Metadata ---
\title{Proposal Title Placeholder}
\author{M. Bronco}
\advisor{Prof. KS}
\department{Department of Electrical and Computer Engineering}
\institution{Boise State University}
\proposaldate{December 2026}

\begin{document}

% =========================================================================
% COVER
% =========================================================================
\maketitle

% =========================================================================
% PROJECT SUMMARY
% =========================================================================
\begin{projectsummary}

  \section*{Project Summary}

  \summaryoverview{

    % space above is needed to allow paragraph below
    Provide a concise overview of the research problem, state of knowledge, central gap,
    hypothesis or design claim, objectives, and general methodology.
  }

  \summaryintellectualmerit{

    % space above is needed to allow paragraph below
    Detail the primary intellectual contributions of the proposed dissertation research,
    distinguishing novel theoretical or engineering findings from standard deliverables.
  }

  \summarybroadersignificance{

    % space above is needed to allow paragraph below
    Explain why resolving this problem matters within ECE and its broader societal,
    technological, or economic impacts.
  }

\end{projectsummary}


% =========================================================================
% PROJECT DESCRIPTION (10–12 Pages Target)
% =========================================================================
\projectdescription

\section{Motivation and Significance}
Define the consequential scientific or technological problem and the context in which it occurs.
Explain why resolving the problem matters and establish the boundaries of the proposed dissertation.

Based on results by Bronco \textit{et al.} \cite{ece_2026_publication},... % cite at least once to prevent build errors.

\section{State of Knowledge and Research Gap}
Synthesize the most relevant literature by claims, mechanisms, approaches, or limitations.
Establish what is known, identify where evidence is incomplete or conflicting, and define the specific
unresolved gap addressed by the dissertation.

\section{Central Question, Hypothesis/Design Claim, and Objectives}
State the organizing research question or engineering challenge, central hypothesis or design claim,
and two to four coherent research objectives.
Explain how completing the objectives will address the central gap.

\begin{objectiveslist}
  \item \textbf{Develop...}
  \item \textbf{Validate...}
  \item \textbf{Implement...}
\end{objectiveslist}

\section{Preliminary Work and Readiness}
Present preliminary evidence that establishes feasibility, validates a method, motivates a design decision, or
refines the proposed direction. Identify limitations and distinguish the student's contributions from those of collaborators.

\section{Research Plan}
For each objective, present the rationale, research question or claim, methods, variables and controls or baselines,
analysis, uncertainty, success criteria, expected outcomes, risks, alternatives, and objective-level deliverables.

\subsection{Objective 1 Methodology}
Detail research activities and experimental design...

\subsection{Objective 2 Methodology}
Detail frameworks and theoretical modeling...

\section{Integration and Contributions}
Explain how the objectives collectively address the central gap. Distinguish research products—such as devices,
algorithms, datasets, models, and publications—from the dissertation's enduring intellectual contribution.

\section{Timeline, Risks, and Completion Plan}
Present a two- to three-year timeline containing research milestones, dependencies, decision points, analysis,
writing, and dissemination. Identify consequential risks, observable triggers, mitigation strategies, and alternative paths.

\begin{riskregister}
  High edge latency & Processing time > 15ms & Optimize matrix kernels in FPGA & Fallback to lower-order estimation \\
  Data dropouts & Packet loss > 5\% & Implement local Kalman prediction & Use historical state extrapolation \\
\end{riskregister}

% =========================================================================
% REFERENCES CITED
% =========================================================================
\newpage
\bibliographystyle{IEEEtran}
\bibliography{ProposalExam}

\end{document}